%% kapa: exam=1 round=31 year=2020 period=1 form=A subject=civil_law subject_ko=민법 seq=5 number=5
%% source: Q-Net 감정평가사 기출문제 (cst00309) artlSeq=5207219 file=extracted/1008/5207219/A형/A형/2020년도 제31회 감정평가사 1차 시험_1교시_A형.pdf
%% conversion: kordoc 4.24.51 -> tools/gen1.py (auto); review=proofread-sonnet; flags=none
\begin{question}{5}
법인에 관한 설명으로 옳지 않은 것은? (다툼이 있으면 판례에 따름)
\begin{choices}
  \item 법인의 대표기관이 법인을 위하여 계약을 체결한 경우, 다른 사정이 없으면 그 성립의 효과는 직접 법인에 미치고 계약을 위반한 때에는 법인이 손해를 배상할 책임이 있다.
  \item 단체의 실체를 갖추어 법인 아닌 사단으로 성립하기 전에 설립주체인 개인이 취득한 권리ㆍ의무는 바로 법인 아닌 사단에 귀속된다.
  \item 법인 아닌 사단은 대표권제한을 등기할 수 없으므로 거래상대방이 사원총회가 대표권제한을 결의한 사실을 몰랐고 모른데 잘못이 없으면, 제한을 넘는 이사의 거래행위는 유효하다.
  \item 민법에서 법인과 그 기관인 이사의 관계는 위임인과 수임인의 법률관계와 같다.
  \item 사단법인의 하부조직 중 하나라 하더라도 스스로 단체의 실체를 갖추고 독자활동을 한다면 독립된 법인 아닌 사단으로 볼 수 있다.
\end{choices}
\end{question}
\answer{2}
